\documentclass{article}
\usepackage{amsmath}
\begin{document}
\section*{Correction: compose the native ray and motion certificates by physical gate}

PR 877 used the index of the closed gate inferred by a qualified
native callback as the index of the positive motion plan. The two
indices have opposite meanings. The main path and timing certificates
in PRs 874 and 876 are unchanged, but their joins to the receiver
crash experiment in PR 877 did not certify collision-free motion.
The point-sampled radio and database outputs of PR 877 are not a
substitute for that missing guarantee.

The repair does not hard-code the opposite label. For each fixed
analytic beam, find the unique closed slab which intersects the exact
six-metre segment. For each of the two motion paths, test the entire
polyline against this concrete slab, both for the center and for
the padded body. The table determines which certified path requires
each physical gate to remain open. On these inputs the first beam
corresponds to the slab crossed by chord 4 of path A, and the second
to the slab crossed by chord 5 of path B (zero-based indices).
The corresponding opposite path has no padded-body intersection.
Both possible orderings of the motion records reconstruct the same
physical choice. The authorization rule now receives the gate index
from the source and selects the certified path for the other, open
gate by this validated relation, rather than reusing an index from
another object. The runtime still uses the explicitly declared
exactly-one-closed premise; two simultaneously closed gates are
not admitted by this correction.

The exact physical join, mesh and motion certificates, analytic
source of the two wire values, native broker, and both real receiver
exit cuts are exercised in an isolated replay. Their previous
aggregate timing and source/radio/settlement counts are reproduced
with the corrected motion assignment. The longer arm egresses
by 453.847823 s, versus the late-evidence lower bound 456.4981695 s
in the declared operating region. The independently selected case
and the limitation about preauthorizing the dock and sharing the
journal remain unchanged.

This is a small example of a concrete comparison obligation:
two internally valid objects with isomorphic finite index sets do
not necessarily agree on their common boundary. The failed
comparison identified two different native effects. This check
establishes the physical relation for the declared contexts only,
not a general Contextual Composition application or a result over
unseen fault mechanisms.
\end{document}
